\documentclass[11pt]{article}
\usepackage[T1]{fontenc}
\usepackage[utf8]{inputenc}
\usepackage{lmodern,microtype,amsmath,amssymb,amsthm,mathtools,booktabs}
\usepackage[margin=27mm]{geometry}
\usepackage[round,authoryear]{natbib}
\usepackage[colorlinks=true,allcolors=blue,breaklinks=true]{hyperref}
\usepackage{url}
\newcommand{\workbenchurl}{https://github.com/KokunoYumeto/schemes-and-moduli-workbench}
\newcommand{\recordurl}{https://github.com/KokunoYumeto/mathematics-workbenches}
\newcommand{\Z}{\mathbb Z}
\newcommand{\Q}{\mathbb Q}
\newcommand{\R}{\mathbb R}
\newcommand{\C}{\mathbb C}
\newcommand{\N}{\mathbb N}
\newcommand{\K}{\mathcal K}
\newcommand{\coker}{\operatorname{coker}}
\newcommand{\ord}{\operatorname{ord}}
\newcommand{\rank}{\operatorname{rank}}
\newcommand{\sgn}{\operatorname{sgn}}
\newcommand{\lcm}{\operatorname{lcm}}
\newcommand{\one}{\mathbf 1}
\newtheorem{theorem}{Theorem}[section]
\newtheorem{proposition}[theorem]{Proposition}
\newtheorem{lemma}[theorem]{Lemma}
\newtheorem{corollary}[theorem]{Corollary}
\theoremstyle{definition}
\newtheorem{definition}[theorem]{Definition}
\newtheorem{example}[theorem]{Example}
\theoremstyle{remark}
\newtheorem{remark}[theorem]{Remark}
\title{Boundary $K$-theory of semilocal adele spaces,\\
after Connes and Consani}
\author{}
\date{24 November 2026}
\begin{document}
\maketitle
\begin{abstract}
Connes and Consani computed the $K$-groups of a semilocal adele algebra with two finite primes after quotienting by compact units. We calculate the boundary map before that quotient, for any finite set of primes. Finite-place invariant groups have finite rank, and their coinvariants are explicit localizations of the integers. A finite integer matrix then determines the kernel and cokernel of the boundary map; for the primes $2,3,5$ it gives a rank-seven kernel and torsion $\Z/2\oplus\Z/4$. We also construct coherent embeddings of a stable place-balancing extension into these semilocal algebras. Its two real character classes become equal in the target. An explicit two-term presentation determines the full middle $K_1$ group, including the extension that the six-term sequence alone leaves unspecified. The mathematics and text were produced by GPT-6 Astra Pro, with separate statement-first re-proofs by the same model and computational checks; no human has checked them.
\end{abstract}

\section{The boundary problem and its answer}

The purpose of this paper is to calculate how compact unit fibres change the boundary $K$-theory of a semilocal adele space. Connes and Consani's calculation after quotienting by those units counts places and their periodic orbits. Before the quotient, two real character classes become equal, and nonconstant unit classes introduce further arithmetic, including torsion. This is a question about a defined crossed-product algebra; it does not require a conjecture about zeros of an $L$-function.

Let $P$ be a nonempty finite set of rational primes. Put
\[
 A_P=\Bigl(\prod_{q\in P}\Q_q\Bigr)\times\R,
 \qquad
 \Gamma_P=\left\{\varepsilon\prod_{q\in P}q^{n_q}:
 \varepsilon\in\{\pm1\},\ n_q\in\Z\right\}.
\]
The discrete group $\Gamma_P$ acts on $A_P$ by diagonal multiplication. Write $W_\varnothing$ for the locus with no zero coordinate, $W_v$ for the locus with exactly the $v$-coordinate zero, and
\[
 W_P^{(1)}=W_\varnothing\cup\bigcup_{v\in P\cup\{\infty\}}W_v.
\]
Define
\[
 N_P=C_0(W_P^{(1)})\rtimes\Gamma_P,\quad
 I_P=C_0(W_\varnothing)\rtimes\Gamma_P,\quad
 D_{v,P}=C_0(W_v)\rtimes\Gamma_P.
\]
All crossed products are full; amenability makes them equal to the reduced crossed products.

Here is the notation needed to state the calculation. Set
\[
 K_P=\prod_{q\in P}\Z_q^\times,
 \qquad
 Y_p=\prod_{q\in P\setminus\{p\}}\Z_q^\times,
 \qquad H_p=\overline{\langle(p)_{q\ne p}\rangle}\subseteq Y_p,
 \qquad \Omega_p=Y_p/H_p.
\]
An empty product is one point. If $U$ runs over the open subgroups of $Y_p$, let $d_U$ be the order of $p$ in $Y_p/U$, and let
\[
 D_p^{\rm rat}=\bigcup_U\frac1{d_U}\Z\subseteq\Q.
\]
This union is directed. Choose a clopen transversal $T_P$ for $c\mapsto-c$ on $K_P$, and let $r:K_P\to T_P$ choose the representative. For $h\in\Z^{\Omega_p}$ let $i_p h$ be its pullback to $K_P$ after dropping the $p$-coordinate. Define
\begin{align}
 \delta_P:\Bigl(\bigoplus_{p\in P}\Z^{\Omega_p}\Bigr)
       \oplus C(T_P,\Z)&\longrightarrow C(K_P,\Z),\label{eq:delta}\\
 \delta_P((h_p),h_\infty)&=-\sum_{p\in P}i_ph_p-r^*h_\infty.\notag
\end{align}

\begin{theorem}[Semilocal boundary groups]\label{thm:main}
For every nonempty finite $P$, each $H_p$ is open, so $\Omega_p$ is finite. There is an exact extension
\[
 0\longrightarrow I_P\longrightarrow N_P
 \longrightarrow \bigoplus_{v\in P\cup\{\infty\}}D_{v,P}\longrightarrow0.
\]
With suspension oriented by increasing logarithmic absolute value,
\[
 K_0(I_P)=0,\qquad K_1(I_P)=C(K_P,\Z),
\]
and the boundary groups are
\begin{align*}
 K_0(D_{p,P})&=\Z^{\Omega_p},&
 K_1(D_{p,P})&=(D_p^{\rm rat})^{\Omega_p},\\
 K_0(D_{\infty,P})&=C(T_P,\Z),&
 K_1(D_{\infty,P})&=0.
\end{align*}
The connecting map is \eqref{eq:delta}. Consequently
\[
 K_0(N_P)=\ker\delta_P
\]
and there is a short exact sequence
\begin{equation}\label{eq:mainK1}
 0\longrightarrow\coker\delta_P\longrightarrow K_1(N_P)
 \longrightarrow\bigoplus_{p\in P}(D_p^{\rm rat})^{\Omega_p}
 \longrightarrow0.
\end{equation}
No splitting of this sequence is part of the assertion.
\end{theorem}

The exact sequence alone does not identify its middle group from the two endpoints. The following presentation determines that extension as well. Set $M_p=C(Y_p,\Z)$ and
\[
 (\sigma_pf)(b)=f(p^{-1}b),\qquad
 (a_pf)(c)=f(\operatorname{pr}_p(c))-f(\operatorname{pr}_p(-c)),\quad c\in T_P.
\]
Define the homomorphism between countable free abelian groups
\begin{align}
 B_P:\bigoplus_{p\in P}M_p&\longrightarrow
 C(T_P,\Z)\oplus\bigoplus_{p\in P}M_p,\label{eq:fullB}\\
 B_P((g_p))&=\left(\sum_pa_pg_p,\ ((1-\sigma_p)g_p)_p\right).\notag
\end{align}

\begin{theorem}[Presentation of the full middle groups]\label{thm:presentation}
For every nonempty finite $P$,
\[
 K_0(N_P)\cong\ker B_P,\qquad K_1(N_P)\cong\coker B_P.
\]
More explicitly, put $L=C(K_P,\Z)$, let $E\subset L$ be its sign-even subgroup, and let $j_p:M_p\to L$ be pullback after dropping the $p$-coordinate. Then
\begin{equation}\label{eq:K1presentation}
 K_1(N_P)\cong
 \frac{L\oplus\bigoplus_pM_p}
 {\left\langle(E,0),\ (j_pg,(1-\sigma_p)g):p\in P,\ g\in M_p\right\rangle}.
\end{equation}
In each relation the second coordinate lies in the $p$th summand. On any cofinal sequence of finite congruence quotients, $B_P$ is the direct limit of the corresponding finite integer matrices, with transition maps given by pullback of functions. Thus this presentation specifies the group extension in \eqref{eq:mainK1}, without assuming a splitting.
\end{theorem}

The kernel and the first subquotient in \eqref{eq:mainK1} admit the smaller finite calculation that follows.

Choose a finite quotient $K_f$ of $K_P$ through which every map to $\Omega_p$ factors and on which $c\mapsto-c$ is free. Choose a transversal $T_f$ and take $T_P$ to be its inverse image. Form the integer matrix
\begin{equation}\label{eq:matrixD}
 D[c,(p,\omega)]=
 \one_{\{\operatorname{pr}_p(c)H_p=\omega\}}
 -\one_{\{\operatorname{pr}_p(-c)H_p=\omega\}},
 \quad c\in T_f,
\end{equation}
where the coset maps are understood on the chosen finite quotient.

\begin{theorem}[Finite matrix and three-prime torsion]\label{thm:finite}
For the matrix \eqref{eq:matrixD},
\[
 K_0(N_P)\cong\ker D,
 \qquad
 \coker\delta_P\cong\coker D\oplus\Z^{(\N)}.
\]
The second isomorphism is not canonical. For $P=\{2,3,5\}$ one may take
\[
 K_f=(\Z/16\Z)^\times\times(\Z/3\Z)^\times
                  \times(\Z/5\Z)^\times.
\]
The matrix has size $32\times14$, rank seven, and nonzero Smith factors
\[
 1,1,1,1,1,2,4.
\]
In particular,
\[
 K_0(N_{\{2,3,5\}})\cong\Z^7,
 \qquad
 \coker\delta_{\{2,3,5\}}
 \cong\Z^{(\N)}\oplus\Z/2\oplus\Z/4.
\]
The torsion subgroup of $K_1(N_{\{2,3,5\}})$ is also $\Z/2\oplus\Z/4$.
\end{theorem}

The last assertion follows from \eqref{eq:mainK1}: its quotient is torsion-free, so every torsion element lies in the left term. These conclusions concern the topology and arithmetic of the semilocal compactification. They neither imply nor refute the Riemann hypothesis.

\paragraph{The place-balancing problem.}
The project's unpublished notes asked whether the row
\[
 (k_p,n_+,n_-)\longmapsto-\sum_{p\in P}k_p-n_+-n_-
\]
is the connecting map of one extension with ideal $C_0(\R)\otimes\K$, with one finite-place quotient for each $p$ and two separate real character classes. The requested comparison with the direct sum of local extensions is a $*$-homomorphism whose ideal map induces addition on $K_1$; the row alone does not construct it. A further requirement is compatibility when primes are added. The notes subsequently gave a stable construction. We prove that construction with explicit local models and realize it inside the semilocal algebra as well.

For the coherent stabilization, let $L_P$ be the multiplicative group of positive rational numbers supported on primes outside $P$, fix $H_0=\ell^2(\N)$, and put
\[
 \widetilde N_P=N_P\otimes\K(\ell^2(L_P)\otimes H_0),\qquad
 J=C_0(\R)\otimes\K(\ell^2(\Q^\times)\otimes H_0).
\]
The generic ideal of $\widetilde N_P$ has the specific identification $C(K_P)\otimes J$ given by the orbit coordinates below. For $P=\varnothing$, use $N_\varnothing=C_0(\R)\rtimes\{\pm1\}$ and $K_\varnothing=\{1\}$.

\begin{theorem}[Coherent place balancing]\label{thm:stable}
For every finite $P$, including the empty set, there is a stable extension
\[
 0\longrightarrow J\longrightarrow E_P^{\rm bal}
 \longrightarrow\left(\bigoplus_{p\in P}\widehat B_p\right)
                  \oplus\widehat B_\infty\longrightarrow0,
\]
where $\widehat B_p\cong C(S^1)\otimes\K$ and
$\widehat B_\infty\cong\C^2\otimes\K$. Its connecting map, in the finite unit and real character coordinates, is the negative-sum row above. There is an injective extension morphism
\[
 E_P^{\rm bal}\longrightarrow\widetilde N_P
\]
whose ideal map is $b\mapsto\one_{K_P}\otimes b$.
For $P\subseteq P'$, the source extensions admit injective morphisms which are the identity on $J$ and insert zero in the new quotient coordinates. The target extensions admit injective morphisms which on generic ideals are
\[
 C(K_P)\otimes J\longrightarrow C(K_{P'})\otimes J,
 \qquad f\otimes b\longmapsto(f\circ\operatorname{pr}_P)\otimes b.
\]
All these morphisms compose for triples of prime sets, and the source-to-target squares commute. Furthermore
\[
 K_0(E_P^{\rm bal})\cong\Z^{|P|+1},\qquad
 K_1(E_P^{\rm bal})\cong\Z^{|P|}.
\]
With the Euclidean pairing of its quotient coordinates, the first group is the root lattice $A_{|P|+1}$.
\end{theorem}

The statement specifies a concrete choice of the local transports used to form the stable extension. It does not identify arbitrary extensions having the same boundary row. The two real labels are distinct in its quotient; their images in the natural boundary have equal $K_0$ class.

\paragraph{Sources and relation to inverse semirings.}
The immediate source is \citet[\S5.3, Proposition 5.4, pp.~23--24; \S5.4, Lemmas 5.5--5.6, pp.~25--28]{CC}, who describe the finite-place mapping tori and compute $K_0=\Z^3$, $K_1=\Z^2$ for the two-prime algebra after quotienting by compact units. An earlier local example appears in \citet[\S10.2, equations (10.10)--(10.19)]{CCM}. Here the full boundary map, its finite orbit arithmetic, and the displayed torsion concern the algebra before the compact-unit quotient. Green imprimitivity, the six-term sequence, and stable extension theory carry the operator-algebra argument.

The definitions use ordinary zeros of local fields. They are not functorially determined by adjoining a new semiring zero. Faul, Goswami, Joubert and Manuell's zero adjunction \citep[Definition 3.6]{FGJM} and inverse-additive-monoid theorem for semimodules \citep[Proposition 6.1]{FGJM} give a different algebraic setting. A support pattern from that setting motivates the stratification, but no equivalence between its semimodules and the $C^*$-algebras above is asserted.

The project's unpublished notes introduced a scalar ``primitive defect'' and developed the semilocal calculation. We write its standard mathematical content as a connecting homomorphism, separate the stable scalar extension from the semilocal extension, and calculate the finite orbit indices and their torsion contribution. The literature search did not establish priority for the assembled formulas; no priority claim is made.

\section{Proof overview}

The main difficulty is that a unit fibre can be harmless for a constant boundary class and decisive for the rest of $K$-theory. We keep those fibres through every orbit calculation. Each individual edge has a free proper action. Its generic quotient is $K_P\times\R$; a finite boundary is a suspension of a profinite rotation; the real boundary is the free sign quotient of $K_P$. The exponential formula for the connecting map then gives the same negative winding on each relevant unit fibre.

The union of the edges has a nonproper action, so its crossed product cannot be replaced by the commutative algebra of its orbit space. Instead, the crossed-product extension and the inclusions of the individual edges assemble the connecting maps by naturality. For the full middle $K_1$, a function on a finite-place transversal determines a unitary on its mapping torus. We lift that unitary over the edge by choosing its real logarithm along the generic line. A coboundary does not lift to zero: it contributes a specified generic winding. These relations give \eqref{eq:K1presentation}. Profinite rotation coinvariants are then computed at finite quotients by averaging point masses on their orbits. For the rational prime rotations considered here, the orbit quotients eventually stabilize: the order at a $q$-power level grows by $q$ at every sufficiently high step. This gives finite-rank invariants and localized-integer coinvariants.

Finally the real boundary fills the sign-even functions. Passing to odd differences leaves a finite matrix. For three primes a seven-column row-lattice calculation gives the torsion directly. The compact-unit quotient supplies a geometric scalar extension and a proper equivariant map into the natural one. Complementary-prime matrix coordinates make these maps compatible as primes are added. Applying fixed orthogonal multiplier corners to the resulting local maps constructs the stable codiagonal and its embedding, including the multiplier identity that equality of boundary rows alone would not give.

\section{Orbit spaces and the extension}

This section supplies the topology needed to apply imprimitivity to the edges, while keeping the whole crossed product intact.

\begin{lemma}[Generic and finite-edge coordinates]\label{lem:orbits}
The action on $W_\varnothing$ is free and proper, with quotient $K_P\times\R$. If $p\in P$, the action on $W_\varnothing\cup W_p$ is also free and proper. Its quotient has generic stratum $K_P\times\R$ and boundary
\[
 B_p=(Y_p\times\R)/((b,s)\sim(pb,s+\log p)).
\]
A generic net $(a_j,t_j)$ converges to $[b,s]\in B_p$ exactly when
\[
 t_j\longrightarrow-\infty,
 \qquad [\operatorname{pr}_p(a_j),t_j]\longrightarrow[b,s]\text{ in }B_p.
\]
\end{lemma}
\begin{proof}
Write a generic point as $((x_q),y)$, put $m_q=v_q(x_q)$ and $u_q=q^{-m_q}x_q$, and let $\varepsilon=\sgn(y)$. Multiplication by the unique element $\varepsilon\prod q^{-m_q}$ makes every finite valuation zero and the real coordinate positive. The resulting coordinates are
\begin{equation}\label{eq:genericcoords}
 a_q=\varepsilon\prod_{r\ne q}r^{-m_r}u_q,
 \qquad t=\log|y|-\sum_{r\in P}m_r\log r.
\end{equation}
They are continuous on each valuation component, invariant under $\Gamma_P$, and have the continuous inverse section $((a_q),e^t)$. This proves the generic quotient assertion.

On $W_p$, normalize the sign and all valuations other than $p$. The unspent power of $p$ acts by $(b,s)\mapsto(pb,s+\log p)$, giving $B_p$. Near that boundary the coordinates before normalizing the $p$-valuation satisfy
\[
 b_q=p^{m_p}a_q\quad(q\ne p),\qquad s=t+m_p\log p.
\]
Convergence of $x_p$ to zero is $m_p\to+\infty$. A representative of the boundary quotient in a local compact section has bounded $s$, and hence $t\to-\infty$. Conversely, the stated two conditions let us choose integers $m_j$ for which $p^{m_j}\operatorname{pr}_p(a_j)\to b$ and $t_j+m_j\log p\to s$. Then $m_j\to+\infty$ and the representatives with $p$-coordinate $p^{m_j}a_{j,p}$ converge to the desired boundary point. No limit of $a_{j,p}$ is necessary, since it is a unit.

Freeness follows from the nonzero real coordinate. To prove properness, let compact sets $C,D$ in this edge satisfy $\gamma C\cap D\ne\varnothing$. For $q\ne p$, the valuations of the nonzero $q$-coordinates on $C,D$ bound the exponent of $q$ in $\gamma$. The real coordinate is bounded away from zero and infinity on both sets, so it bounds $\log|\gamma|$. The remaining exponent of $p$ is therefore bounded too. The sign has two possibilities. Only finitely many $\gamma$ occur, which is the properness criterion for a discrete action.
\end{proof}

\begin{lemma}[Real edge]\label{lem:real}
After dividing by the positive prime-power directions, the real edge is $K_P\times\R$ with the residual action $(c,z)\mapsto(-c,-z)$. This action is free. Its quotient is $T_P\times\R$, and its boundary is $T_P$.
\end{lemma}
\begin{proof}
All finite coordinates are nonzero on this edge. Their valuations normalize them to units, leaving only the common sign. For properness those same valuations bound every prime exponent between compact sets. The sign action is free because $c=-c$ would imply $2c_q=0$ in a characteristic-zero field.

To construct $T_P$, use one odd prime coordinate modulo that prime and choose one residue from each pair $\{a,-a\}$; if only the prime $2$ is available, use residues modulo $4$. Pullback gives a clopen transversal. Write $c=\varepsilon h$ with $h\in T_P$. The map $[(\varepsilon h,z)]\mapsto(h,\varepsilon z)$ is a homeomorphism to $T_P\times\R$. On its two generic rays, if the second coordinate is $w\ne0$, its relation to \eqref{eq:genericcoords} is
\[
 a=\sgn(w)h,\qquad t=\log|w|.
\]
\end{proof}

\begin{proposition}[Extension and failure of global properness]\label{prop:extension}
The extension in Theorem~\ref{thm:main} is exact. The action on $W_P^{(1)}$ is not proper.
\end{proposition}
\begin{proof}
The complement of $W_P^{(1)}$ is the finite union of the intersections of two coordinate hyperplanes, hence is closed. Within $W_P^{(1)}$, the generic stratum is open and its complement is the finite disjoint union of the closed strata $W_v$. Restriction of continuous functions therefore gives an exact equivariant sequence. Full crossed products preserve exact sequences by their universal property; amenability of $\Gamma_P$ identifies this with the reduced sequence.

For nonproperness fix $p\in P$. Let $x_n$ have $p$-coordinate $p^n$, all other finite coordinates $1$, and real coordinate $1$. Then $x_n$ converges to a point of $W_p$. The sequence $p^{-n}x_n$ has $p$-coordinate $1$ and real coordinate $p^{-n}\to0$. Its other finite coordinates lie in compact unit groups. A subsequence converges to a point of $W_\infty$. The source and target sequences thus lie in compact subsets of $W_P^{(1)}$, whereas the acting elements $p^{-n}$ leave every finite subset of $\Gamma_P$. This contradicts properness.
\end{proof}

For a free proper action, Green imprimitivity identifies the crossed product with the commutative quotient algebra up to strong Morita equivalence \citep[Theorem 14 and Corollary 15, pp.~88--91]{Green}. Its construction restricts to an invariant open subset and descends to its closed complement: in the completion of $C_c(X)$, the submodule generated by $C_0(U/G)$ is the completion of $C_c(U)$, and the quotient is the corresponding module for the closed complement. Taking compact operators therefore gives compatible equivalences for the three terms of an edge extension. We use this only for the proper actions proved above.

\section{The boundary maps with unit fibres}

This section computes the connecting homomorphism on every boundary class. Throughout, a loop with argument increasing by $2\pi$ along the real suspension coordinate represents $+1$.

\begin{lemma}[Profinite coefficients and suspensions]\label{lem:profinite}
For a compact metrizable zero-dimensional space $Y$,
\[
 K_0(C(Y))=C(Y,\Z),\qquad K_1(C(Y))=0.
\]
If $\sigma$ is induced by a homeomorphism of $Y$, then its mapping torus has
\[
 K_0(M_\sigma)=\ker(1-\sigma_*),\qquad
 K_1(M_\sigma)=\coker(1-\sigma_*)
\]
on $C(Y,\Z)$.
\end{lemma}
\begin{proof}
Finite clopen partitions give finite-dimensional subalgebras of $C(Y)$ whose union is dense. On each partition, a projection is classified stably by its integer rank on each piece, and the odd group is zero. Continuity of $K$-theory gives the first two statements. Evaluation at an endpoint of the mapping torus has suspended kernel. Its connecting map is $1-\sigma_*$ under Bott identification \citep[Proposition 10.4.1, p.~76]{Blackadar}. The six-term sequence and the first two statements give the asserted kernel and cokernel.
\end{proof}

\begin{proposition}[Finite-place boundary]\label{prop:finiteboundary}
Let $F_p=C(Y_p/H_p,\Z)$ and $C_p=C(Y_p,\Z)/(1-\sigma_p)C(Y_p,\Z)$, where $(\sigma_pf)(y)=f(p^{-1}y)$. Then
\[
 K_0(D_{p,P})=F_p,\qquad K_1(D_{p,P})=C_p.
\]
In these coordinates the connecting map of the $p$-edge is $h\mapsto-i_ph$.
\end{proposition}
\begin{proof}
Lemma~\ref{lem:orbits} identifies the boundary quotient with $B_p$. Functions on this suspension are the mapping torus of $\sigma_p$ after rescaling its time by $\log p$. A locally constant function fixed by $p$ is fixed by the closure $H_p$, so Lemma~\ref{lem:profinite} gives the two groups.

It remains to determine the map and its sign. Let $E$ be clopen in $Y_p/H_p$ and let $\widetilde E$ be its inverse image in $Y_p$. On the edge quotient define
\[
 x_E(a,t)=\frac{\one_{\widetilde E}(\operatorname{pr}_p(a))}{1+e^t},
 \qquad x_E([b,s])=\one_{\widetilde E}(b).
\]
The boundary expression is well defined because $\widetilde E$ is $p$-invariant. Lemma~\ref{lem:orbits} proves continuity across the boundary. The function vanishes at the remaining end $t\to+\infty$, and its sets above a positive threshold are compact because the unit fibre and the suspension boundary are compact. It is therefore a self-adjoint lift of the boundary projection.

The exponential formula for the connecting map \citep[\S9.3.2, p.~68]{Blackadar} gives $\exp(2\pi ix_E)$. On every fibre above $\widetilde E$ its argument decreases from $2\pi$ to $0$ as $t$ increases. Its winding is $-1$ there, and zero elsewhere. Characteristic functions of clopen sets generate $F_p$, which proves the formula on every class. The projection $K_P\to Y_p/H_p$ is surjective, so this boundary map is injective.
\end{proof}

\begin{proposition}[Real-place boundary and character collapse]\label{prop:realboundary}
The real boundary has $K_0=C(T_P,\Z)$ and $K_1=0$. Its connecting map is $h\mapsto-r^*h$. Under the crossed-product coordinate $C(K_P)\rtimes C_2$, the two projections $e_\pm=(1\pm u)/2$ have the same $K_0$ class, and each maps to $-\one_{K_P}$.
\end{proposition}
\begin{proof}
The clopen transversal gives
\[
 C(K_P)\rtimes C_2\cong M_2(C(T_P)).
\]
The projections $e_+,e_-$ become two constant rank-one projections, conjugate by a constant unitary. Thus their classes coincide. The group computation follows from Lemma~\ref{lem:profinite}.

In the orbit coordinate $(h,w)\in T_P\times\R$ of Lemma~\ref{lem:real}, a lift of a clopen projection $\one_E$ at $w=0$ is $\one_E(h)/(1+|w|)$. On both rays, increasing $t=\log|w|$ makes its exponential argument decrease from $2\pi$ to zero. The winding is $-1$ on $E\cup(-E)$ in $K_P$. By additivity the map is $-r^*$, whose image is exactly the subgroup of sign-even functions.
\end{proof}

\begin{proof}[Boundary-map part of Theorem~\ref{thm:main}]
The generic quotient of Lemma~\ref{lem:orbits} gives $K_0(I_P)=0$ and $K_1(I_P)=C(K_P,\Z)$ by suspension. Inclusion of each invariant open edge into $W_P^{(1)}$ induces a morphism of crossed-product extensions, identical on the generic ideal and the named summand inclusion on the quotient. Naturality and additivity of connecting maps therefore sum the formulas of Propositions~\ref{prop:finiteboundary} and \ref{prop:realboundary}. This is \eqref{eq:delta}. Six-term exactness \citep[Theorem 9.3.1, pp.~67--68]{Blackadar} then gives $K_0(N_P)=\ker\delta_P$ and \eqref{eq:mainK1}, with $C_p$ temporarily in place of $(D_p^{\rm rat})^{\Omega_p}$. Section~\ref{sec:arithmetic} proves that identification and the finiteness of $\Omega_p$.
\end{proof}

\begin{proof}[Proof of Theorem~\ref{thm:presentation}]
Fix $p$ and put $\ell=\log p$. Let $X_p$ be the proper $p$-edge quotient of Lemma~\ref{lem:orbits}. For $f\in M_p$, define a unitary on its boundary by
\[
 u_f([b,s])=\exp(2\pi i(s/\ell)f(b)),\qquad0\leq s\leq\ell.
\]
The two endpoint values are both $1$, so this is well defined on the mapping torus. Inclusion of the suspended coefficient algebra in that torus sends $f$ to $[u_f]$. Lemma~\ref{lem:profinite} therefore identifies this class with the image of $f$ in $M_p/(1-\sigma_p)M_p$.

On $Y_p\times\R$, construct a continuous real logarithm $\phi_f$ by requiring
\[
 \phi_f(b,0)=0,\qquad
 \exp(2\pi i\phi_f(b,t))=u_f([b,t]).
\]
For an explicit construction, on every interval $[k\ell,(k+1)\ell]$ give $\phi_f(b,t)$ slope $f(p^{-k}b)/\ell$ in $t$, and choose successive endpoint values by continuity, starting at $0$. Each bounded time interval uses finitely many locally constant functions of $b$, so the resulting function is continuous in both variables. The construction is additive in $f$.

Choose a continuous function $\chi:\R\to[0,1]$ equal to $1$ on $(-\infty,0]$ and to $0$ on $[1,\infty)$. Set
\[
 U_f(a,t)=\exp\bigl(2\pi i\chi(t)\phi_f(\operatorname{pr}_p(a),t)\bigr)
\]
on the generic part of $X_p$, and set $U_f=u_f$ on the boundary. At all sufficiently negative times this is $u_f([\operatorname{pr}_p(a),t])$, so the edge convergence criterion proves continuity at the boundary. At positive large times it is $1$. The closure of $K_P\times(-\infty,1]$ in $X_p$ is compact: after choosing a bounded representative for the mapping-torus time, an escaping negative-time sequence has a convergent boundary subsequence, and the unit groups and boundary torus are compact. Thus $U_f-1\in C_0(X_p)$. Addition of functions multiplies these scalar unitaries, so $f\mapsto[U_f]$ is a group homomorphism. Transfer it by the edge Morita module and then include the edge in $N_P$; denote the resulting homomorphism by
\[
 \kappa_p:M_p\longrightarrow K_1(N_P).
\]
Its boundary restriction is the class of $f$ in the $p$th coinvariant group, with zero in the other summands.

We compute its relation on a coboundary. For $g\in M_p$ define the real function
\[
 G_g([b,s])=(1-s/\ell)g(b)+(s/\ell)g(p^{-1}b),\qquad0\leq s\leq\ell.
\]
Its endpoint equality is precisely the mapping-torus identification, so it is a continuous real function on $B_p$. If $d_p=1-\sigma_p$, then
\[
 u_{d_pg}=\exp(-2\pi iG_g),\qquad
 \phi_{d_pg}(b,t)=g(b)-G_g([b,t]).
\]
The latter identity follows either on the defining intervals or from uniqueness of the continuous real logarithm with value zero at $t=0$.

It follows that $U_{d_pg}$ is the product of the generic unitary
$\exp(2\pi i\chi(t)j_pg(a))$ and the unitary whose generic formula is
\[
 \exp\bigl(-2\pi i\chi(t)G_g([\operatorname{pr}_p(a),t])\bigr).
\]
The second unitary is null-homotopic in the unitization of $C_0(X_p)$: multiply its real exponent by $\lambda\in[0,1]$, with boundary value $\exp(-2\pi i\lambda G_g)$. The function $G_g$ is bounded on the compact boundary torus, so this is a norm-continuous homotopy, including at $\lambda=0$. The first is $1$ on the boundary and at both generic ends; its winding along increasing $t$ is $-j_pg$. Consequently
\begin{equation}\label{eq:coboundarylift}
 \kappa_p((1-\sigma_p)g)=-[j_pg]\quad\hbox{in }K_1(N_P),
\end{equation}
where the bracket on the right means the image of the generic $K_1$ group.

Map $L\oplus\bigoplus_pM_p$ to $K_1(N_P)$ using the generic inclusion on $L$ and the maps $\kappa_p$. It is surjective: it surjects onto all the boundary $K_1$ groups, and six-term exactness says that the remaining kernel of boundary restriction is the generic image. The relations in \eqref{eq:K1presentation} vanish by the real-place boundary formula and \eqref{eq:coboundarylift}.

Conversely, suppose $(l,(f_p))$ maps to zero. Its boundary restriction gives $f_p=(1-\sigma_p)g_p$ for some $g_p\in M_p$. Subtracting the relations $(j_pg_p,(1-\sigma_p)g_p)$ leaves $(l-\sum_pj_pg_p,0)$. The latter generic class lies in $\operatorname{im}\delta_P$, which is
\[
 E+\sum_pj_p\ker(1-\sigma_p).
\]
The first summand is a defining relation. Each element $j_pg$ of the second is also a defining relation, since $(1-\sigma_p)g=0$. There are therefore no further relations, proving \eqref{eq:K1presentation}.

Odd difference identifies $L/E$ with $C(T_P,\Z)$ and sends $j_p$ to $a_p$. This proves the cokernel formula. The kernel of $B_P$ consists exactly of tuples of invariant functions for which $\sum_pj_pg_p$ is sign-even. As in the proof of Theorem~\ref{thm:finite}, the real boundary coordinate is then uniquely determined. This gives $\ker B_P\cong K_0(N_P)$.

For the last assertion, choose nested congruence quotients of $K_P$ that distinguish sign, and use the inverse image of one initial sign transversal at every further level. Their projections give nested finite quotients of each $Y_p$. Every locally constant integer function factors through one of them. Translation, pullback, and odd difference preserve these levels. Thus \eqref{eq:fullB} is the direct limit of its finite integer matrices. Filtered colimits of abelian groups preserve kernels and cokernels, giving the claimed explicit direct-system description of both groups.
\end{proof}

\begin{example}[One prime]\label{ex:oneprime}
If $P=\{p\}$, then $Y_p$ is one point and $F_p=C_p=\Z$. The boundary map is
\[
 (n,h)\longmapsto-n\one_{K_P}-r^*h.
\]
Its kernel consists of $(n,-n\one_{T_P})$, so $K_0(N_P)=\Z$. Odd difference identifies its cokernel with $C(T_P,\Z)$. Hence
\[
 0\longrightarrow C(T_P,\Z)\longrightarrow K_1(N_P)
 \longrightarrow\Z\longrightarrow0.
\]
This sequence splits abstractly because $\Z$ is free. It gives $K_1(N_P)\cong C(T_P,\Z)\oplus\Z$, without choosing a preferred splitting.
\end{example}

\section{Arithmetic orbit indices and rational coinvariants}\label{sec:arithmetic}

Finite-quotient Haar normalization is also used in odometer $K$-theory: \citet[Proposition 4.10, pp.~34--35]{GPW} obtain the trace range from reciprocal finite indices. The calculation below records the full coinvariant group, including the orbit quotient when the rotation is not minimal on the entire profinite group.


This section makes the profinite rotation groups explicit and proves that the invariant part has finite rank.

\begin{theorem}[Haar coordinate for coinvariants]\label{thm:haar}
Let $Y$ be a compact metrizable profinite abelian group, let $g\in Y$, and put $H=\overline{\langle g\rangle}$. For an open subgroup $U\leq Y$, write $d_U$ for the order of $g$ in $Y/U$, and put $D_H=\bigcup_Ud_U^{-1}\Z$. Then normalized Haar averaging induces an isomorphism
\[
 C(Y,\Z)/(1-g)C(Y,\Z)\ \xrightarrow{\ \sim\ }\ C(Y/H,D_H),
\]
where the right-hand coefficient group is discrete. The map is
\[
 [f]\longmapsto\left(yH\longmapsto\int_H f(yh)\,dh\right).
\]
\end{theorem}
\begin{proof}
Every integer-valued continuous function factors through a finite quotient $Y/U$. Translation by $g$ partitions this finite group into orbits, all of length $d_U$. On one orbit the quotient by the differences of adjacent point indicators is free cyclic: all point indicators agree, and summation of their coefficients proves that there is no further relation. Haar averaging sends its point generator to $1/d_U$ on that orbit. Thus averaging is injective on finite-level coinvariants.

If $V\subseteq U$, each orbit in $Y/V$ above a given orbit of $Y/U$ maps onto it with degree $d_V/d_U$. Pullback of one point indicator therefore contributes $d_V/d_U$ point generators to each fine orbit. Their averaged value is
\[
 \frac{d_V}{d_U}\frac1{d_V}=\frac1{d_U},
\]
so the finite-level injections are compatible with averaging. Filtered colimits commute with cokernels, hence the induced map on the displayed quotient is injective.

For surjectivity, a continuous map from compact $Y/H$ to discrete $D_H$ has finite image and a finite clopen level partition. Pull the partition back to $Y$; it is invariant under $H$ and factors through some $Y/U$. Its finitely many rational values have denominators dividing some $d_V$, after taking a common refinement. At a further quotient with both properties, any prescribed value $a/d_V$ is an integral multiple of the average $1/d_V$ of a point generator on the corresponding orbit. Choosing those multiples on all orbits realizes the function. This proves surjectivity and the theorem.
\end{proof}

\begin{theorem}[Finite orbit index and denominator formula]\label{thm:indices}
Fix distinct rational primes $p$ and the members of a finite set $Q$. Put $Y=\prod_{q\in Q}\Z_q^\times$ and $H=\overline{\langle p\rangle}$. If $Q$ is empty, $Y/H$ is one point and $D_H=\Z$. Otherwise define
\[
 d_q=\begin{cases}\ord_q(p),&q\ne2,\\ \ord_4(p),&q=2,\end{cases}
 \quad c_q=v_q(p^{d_q}-1),\quad b_q=v_q(d_q),
\]
and let $d_*$ be the part of $\lcm_{q\in Q}d_q$ supported on primes outside $Q$. Then $H$ is open and
\begin{equation}\label{eq:hformula}
 [Y:H]=\frac{\prod_{q\in Q}(q-1)q^{c_q-b_q-1}}{d_*},
 \qquad
 D_H=\frac1{d_*}\Z\left[\frac1{\prod_{q\in Q}q}\right].
\end{equation}
\end{theorem}
\begin{proof}
For odd $q$, $d_q$ is prime to $q$ and $c_q\geq1$. For $q=2$, $d_q\in\{1,2\}$ and $c_q\geq2$. The integer $p^{d_q}$ is not $1$, so $c_q$ is finite.

If $a\equiv1\pmod q$ with $v_q(a-1)=c\geq1$, and additionally $c\geq2$ for $q=2$, the binomial expansion proves
\[
 v_q(a^q-1)=c+1.
\]
Indeed the first term is $q(a-1)$; every later term has valuation strictly greater than $c+1$ under the stated hypotheses. Induction gives $v_q(a^{q^k}-1)=c+k$. Raising to an exponent prime to $q$ preserves valuation, by factoring the geometric sum, whose value is a unit modulo $q$. It follows that
\begin{equation}\label{eq:orderformula}
 \ord_{q^m}(p)=d_q q^{\max(0,m-c_q)}
\end{equation}
for $m\geq1$ when $q$ is odd and $m\geq2$ when $q=2$.

At finite level $Y_{\boldsymbol m}=\prod_q(\Z/q^{m_q}\Z)^\times$, every $p$-orbit has length
\[
 L_{\boldsymbol m}=\lcm_{q\in Q}\bigl(d_qq^{\max(0,m_q-c_q)}\bigr).
\]
Once every $m_q$ is large enough that $m_q-c_q+b_q$ dominates the $q$-valuation of all $d_r$, this becomes
\[
 L_{\boldsymbol m}=d_*\prod_{q\in Q}q^{m_q-c_q+b_q}.
\]
Divide $|Y_{\boldsymbol m}|=\prod_q(q-1)q^{m_q-1}$ by this expression. The quotient is the fixed integer in \eqref{eq:hformula}. The finite orbit quotients therefore have constant cardinality at all further levels, and their surjective transition maps are bijective. Their inverse limit $Y/H$ is this finite set. Thus $H$ is open.

Finally the denominators $L_{\boldsymbol m}$ have the fixed outside part $d_*$ and arbitrarily large powers of every $q\in Q$. Their reciprocal integer groups have union exactly the displayed subgroup of $\Q$. This proves both formulas.
\end{proof}

Applying Theorems~\ref{thm:haar} and \ref{thm:indices} with $Q=P\setminus\{p\}$ completes Theorem~\ref{thm:main}. If $|P|\geq2$, each $D_p^{\rm rat}$ has unbounded divisibility by a prime in $P\setminus\{p\}$. It is not a nonzero free abelian group. Thus a claim that these coinvariants are merely free groups on the final orbit sets would lose their essential denominator data.

\section{The finite matrix and the three-prime example}

This section proves the finite reduction and gives the torsion calculation in elementary integer coordinates.

\begin{proof}[Finite-reduction part of Theorem~\ref{thm:finite}]
Let $S_h=\sum_pi_ph_p$. Equation \eqref{eq:delta} vanishes exactly when $S_h$ is sign-even; in that case $h_\infty$ is uniquely $-S_h|_{T_P}$. Odd difference
\[
 C(K_P,\Z)\longrightarrow C(T_P,\Z),\qquad
 f\longmapsto\bigl(c\mapsto f(c)-f(-c)\bigr)
\]
is surjective, since a function on $T_P$ can be extended by zero on $-T_P$. Its kernel is precisely the image of the real boundary. Thus $K_0(N_P)$ is the kernel of the finite-place odd-difference map, and $\coker\delta_P$ is its cokernel. All its finite-place functions factor through the chosen $K_f$. This map is exactly the matrix $D$ of \eqref{eq:matrixD} followed by pullback from $T_f$.

Choose an increasing sequence of finite clopen partitions of $T_P$ beginning with $T_f$ and separating points. At each refinement, the inclusion of integer function groups splits: choose one child of each old cell and recover the old value from that child. Its complement is a finite-rank free group of differences. Hence
\[
 C(T_P,\Z)\cong\Z^{T_f}\oplus\Z^{(\N)}.
\]
The free complement has infinite rank because every nonempty clopen subset of the product of local unit groups has arbitrarily fine nontrivial clopen partitions. The image of $D$ lies in the first summand. Taking kernel and cokernel proves the two assertions.
\end{proof}

\begin{proof}[Three-prime computation]
Take $P=\{2,3,5\}$. Formula \eqref{eq:hformula}, or the orders at the displayed finite level, gives
\[
 |\Omega_2|=2,\quad |\Omega_3|=8,\quad |\Omega_5|=4,
\]
and
\[
 D_2^{\rm rat}=\tfrac14\Z[1/15],\quad
 D_3^{\rm rat}=\Z[1/10],\quad
 D_5^{\rm rat}=\Z[1/6].
\]
There are $8\cdot2\cdot4=64$ unit tuples at moduli $16,3,5$, so a sign transversal has $32$ points. The $14$ columns of $D$ come in opposite pairs under $\omega\mapsto-\omega$; no pair is fixed. Delete one column of each pair. This does not change the column image, and leaves a $32\times7$ matrix.

Order its seven columns in blocks of sizes $1,4,2$. Representatives of the orbit pairs, in the remaining finite coordinates, are
\[
 \begin{array}{c|l}
 p&\text{representatives}\\\hline
 2&(1,1)\\
 3&(1,1),(1,2),(1,3),(1,4)\\
 5&(1,1),(1,2).
 \end{array}
\]
Here entries in the $p=3$ row are modulo $16,5$, and those in the $p=5$ row are modulo $16,3$. Let $e_j$ denote the $j$th standard vector in the middle four-column block, and set $(s_1,s_2,s_3,s_4)=(1,-1,-1,1)$. Multiplying these representatives by their corresponding primes modulo the indicated moduli gives, up to overall row sign, exactly the rows
\begin{equation}\label{eq:rows}
 (1,\pm e_j,s_j,0),
 \qquad (1,\pm e_j,0,-s_j),\qquad 1\leq j\leq4.
\end{equation}
This orbit table can also be checked directly: multiplication by $2$ on the $3,5$ coordinates has order four; multiplication by $3$ on the $16,5$ coordinates has order four; multiplication by $5$ on the $16,3$ coordinates has order four. The listed representatives and their negative orbits exhaust respectively $8,32,16$ elements.

Write $E_1,\ldots,E_7$ for the standard basis of $\Z^7$. From \eqref{eq:rows}, the row lattice is generated by
\[
 E_1+E_{j+1}+s_jE_6,\quad 2E_{j+1}\ (1\leq j\leq4),
 \quad E_6+E_7.
\]
Both inclusions are immediate: row differences give the last five generators, and each row differs from the first generator of its index by one of them. In the quotient, eliminate $E_{j+1}=-E_1-s_jE_6$ and $E_7=-E_6$. The remaining relations are
\[
 2(E_1+E_6)=0,\qquad2(E_1-E_6)=0.
\]
Their $2\times2$ relation matrix has entries with gcd $2$ and determinant of absolute value $8$, so the quotient is $\Z/2\oplus\Z/4$. The row lattice consequently has full rank seven and index eight. Transposition preserves the nonzero Smith factors; these are five $1$'s followed by $2,4$. The original matrix has kernel rank $14-7=7$ and cokernel free rank $32-7=25$. The finite reduction absorbs that finite free rank into the countable free summand, proving all assertions. The claim about torsion of $K_1$ follows as explained after Theorem~\ref{thm:finite}.
\end{proof}

An independent finite-group enumeration, the program \emph{Semilocal Boundary Check} in the \href{\workbenchurl}{workbench}, checks the orbit orders, the normalized-Haar transition coefficients, and Smith forms for six prime sets. The argument above supplies the proof of the three-prime calculation.

\section{The compact-unit quotient and coherent inclusions}

This section constructs the algebra maps needed to compare place balancing with the semilocal algebra. It is useful to do this before choosing any stable Busby sum.

Let $Z_P=W_P^{(1)}/K_P$, where $K_P$ acts independently by multiplication on the finite coordinates, and let
\[
 S_P=C_0(Z_P)\rtimes\Gamma_P,
 \qquad
 \widetilde S_P=S_P\otimes\K(\ell^2(L_P)\otimes H_0).
\]
The compact group action commutes with the $\Gamma_P$ action. For $P=\varnothing$ these definitions give the reflection crossed product. Write $\widetilde I_P^{\rm sc}$ for the generic ideal of $\widetilde S_P$.

\begin{lemma}[The geometric scalar comparison]\label{lem:scalarcomparison}
The quotient map $q_P:W_P^{(1)}\to Z_P$ is proper and equivariant. Pullback and crossed product give an injective morphism of the generic-boundary extensions
\[
 \iota_P:\widetilde S_P\longrightarrow\widetilde N_P.
\]
There are specific identifications
\[
 \widetilde I_P^{\rm sc}=J,\qquad
 \widetilde I_P=C(K_P)\otimes J
\]
under which its ideal map is $j_P(b)=\one_{K_P}\otimes b$. Both generic ideals are essential.
\end{lemma}
\begin{proof}
For a compact group acting on a locally compact Hausdorff space, its orbit map is closed and has compact fibres. To see closedness, if a net $k_i x_i$ from the saturation of a closed set converges, pass to a subnet with $k_i$ convergent; then $x_i$ converges to a point of that closed set. The quotient is locally compact Hausdorff, and this closed map with compact fibres is proper. Thus $q_P^*$ maps $C_0(Z_P)$ injectively and nondegenerately into $C_0(W_P^{(1)})$. It respects the open generic stratum and its closed complement. Equivariance is the commutation of the two multiplication actions.

An injective equivariant coefficient map induces an injective reduced crossed-product map: use faithful regular representations, or the faithful coefficient expectations. The groups here are amenable, so the same conclusion holds for full crossed products. Exactness gives the stated extension morphism, including its injective map on the boundary quotient.

For a generic point $x=((x_q),y)$ put
\[
 g(x)=\sgn(y)\prod_{q\in P}q^{v_q(x_q)}\in\Gamma_P.
\]
Then $x=g(x)((a_q),e^t)$ with $(a,t)$ as in \eqref{eq:genericcoords}. This is an equivariant homeomorphism
\[
 W_\varnothing\cong\Gamma_P\times K_P\times\R,
\]
where $\Gamma_P$ translates its own factor and fixes the other two. The corresponding generic part of $Z_P$ is $\Gamma_P\times\R$, and $q_P$ is the projection forgetting $a$. Crossing the translation factor gives
\[
 I_P\cong C(K_P)\otimes C_0(\R)\otimes\K(\ell^2(\Gamma_P)),
\]
with the analogous formula omitting $C(K_P)$ in the scalar case. The fixed basis bijection
\[
 \Gamma_P\times L_P\longrightarrow\Q^\times,
 \qquad(\gamma,l)\longmapsto\gamma l
\]
then gives the asserted stabilized identifications and the exact formula for $j_P$.

Finally, the generic stratum is dense. If $a$ in either crossed product annihilates its generic ideal, then for every generic compactly supported coefficient $f$,
\[
 f^*E(a^*a)f=E(f^*a^*af)=0,
\]
where $E$ is the faithful coefficient expectation. Thus $E(a^*a)$ vanishes on the generic stratum, hence everywhere by continuity. Faithfulness gives $a=0$. Tensoring with compact operators preserves essentiality.
\end{proof}

\begin{proposition}[Coherence when primes are added]\label{prop:coherence}
For all finite $P\subseteq P'$ there are injective extension morphisms
\[
 \Psi^{\rm sc}_{P,P'}:\widetilde S_P\longrightarrow\widetilde S_{P'},
 \qquad
 \Psi^{\rm nat}_{P,P'}:\widetilde N_P\longrightarrow\widetilde N_{P'}.
\]
They compose exactly for triples of prime sets. Their ideal maps are respectively the identity of $J$ and $\operatorname{pr}_P^*\otimes\mathrm{id}_J$. The squares with $\iota_P,\iota_{P'}$ commute. On the boundary, the maps have zero components at the new primes. Their ranges lie in the ideals where every new prime coordinate is nonzero.
\end{proposition}
\begin{proof}
Put $R=P'\setminus P$, let $G_R$ be the positive rational group generated by $R$, and let $U_R\subset W_{P'}^{(1)}$ be the locus where all coordinates indexed by $R$ are nonzero. For $x\in U_R$ set
\[
 h_R(x)=\prod_{r\in R}r^{v_r(x_r)},\qquad
 z_v=x_v/h_R(x)\ (v\in P\cup\{\infty\}),\qquad
 b_r=x_r/h_R(x)\ (r\in R).
\]
Each $b_r$ is an $r$-adic unit. These formulas give a homeomorphism
\begin{equation}\label{eq:oldsector}
 U_R\longrightarrow G_R\times W_P^{(1)}\times K_R,
 \qquad x\longmapsto(h_R,z,b),
\end{equation}
whose inverse multiplies every coordinate of $(z,b)$ by $h_R$. Under $\Gamma_{P'}=\Gamma_P\times G_R$, the element $(\gamma,\rho)$ acts by
\[
 (h,z,b)\longmapsto(\rho h,\gamma z,\gamma b).
\]
In particular the new unit coordinates are acted on by $\Gamma_P$; they cannot be discarded from this action.

The crossed product of the translation factor $G_R$ is $\K(\ell^2(G_R))$. Consequently the crossed product on $U_R$ is
\[
 \K(\ell^2(G_R))\otimes
 \bigl(C_0(W_P^{(1)}\times K_R)\rtimes\Gamma_P\bigr).
\]
The proper equivariant projection to $W_P^{(1)}$ gives an injective map from $N_P$ to the second factor. Combine it with
\[
 L_P=G_R\times L_{P'},
 \qquad(\rho,l)\longmapsto\rho l,
\]
and extend the resulting elements by zero from the invariant open set $U_R$. This defines $\Psi^{\rm nat}_{P,P'}$. In the scalar quotient, \eqref{eq:oldsector} becomes $G_R\times Z_P$; the same construction without $K_R$ defines $\Psi^{\rm sc}_{P,P'}$. Both are injective by the crossed-product factorization just given. They preserve the generic ideals and boundary restrictions, and a new zero-coordinate stratum is disjoint from $U_R$.

For clarity, the formula on a dense algebra determines all choices. If $e_{\rho,\sigma}$ is a matrix unit for $\ell^2(G_R)$, $f\in C_c(W_P^{(1)})$, and $u_\gamma$ is the crossed-product symbol with $u_\gamma f u_\gamma^*=f\circ\gamma^{-1}$, then
\begin{equation}\label{eq:matrixinclusion}
 (f u_\gamma)\otimes e_{\rho,\sigma}
 \longmapsto
 \left[\one_{\{h_R=\rho\}}
 f\!\left(x_{P\cup\{\infty\}}/h_R(x)\right)\right]
 u_{\rho\sigma^{-1}\gamma}.
\end{equation}
The remaining $\ell^2(L_{P'})\otimes H_0$ matrix factor is unchanged. The coefficient in brackets has compact support in $U_R$: $h_R$ is fixed, $f$ has compact support, and $K_R$ is compact. It therefore extends by zero continuously to $W_{P'}^{(1)}$. The scalar formula is identical with $f$ a function on $Z_P$.

Now take $P\subseteq P'\subseteq P''$ and put $T=P''\setminus P'$. Dividing the old coordinates by $h_T$ does not change their $R$-valuations, since a rational number supported on $T$ is an $r$-adic unit for every $r\in R$. Applying \eqref{eq:matrixinclusion} twice therefore gives coefficient
\[
 \one_{\{h_R=\rho,\ h_T=\lambda\}}
 f\!\left(x_{P\cup\{\infty\}}/(h_Rh_T)\right)
\]
and group symbol $u_{\rho\sigma^{-1}\lambda\mu^{-1}\gamma}$. These are exactly the direct formula for $R\cup T$, under the fixed basis multiplication map. Such tensors span a dense subalgebra, proving strict composition. Pulling a scalar coefficient back to $W_P^{(1)}$ before or after this formula gives the same function. This proves the commuting square.

On a generic point, write $x=g(x)(a,e^t)$ as in Lemma~\ref{lem:scalarcomparison}. In \eqref{eq:oldsector}, $g(x)=h_R(x)g(z)$, the old unit coordinate of $z$ is $\operatorname{pr}_P(a)$, and its invariant logarithm is $t$. Matrix units are relabelled by precisely the same multiplication that identifies $\Gamma_P L_P$ and $\Gamma_{P'}L_{P'}$ with $\Q^\times$. Thus \eqref{eq:matrixinclusion} is $\operatorname{pr}_P^*\otimes\mathrm{id}_J$ on the natural generic ideals, and is the identity on the scalar ideals.
\end{proof}

For a proper inclusion of prime sets these middle-algebra maps are nondegenerate into their old-stratum ideals, but need not be nondegenerate into the entire target algebra. Their generic-ideal maps are nondegenerate. The latter fact is the one needed to extend those maps to multiplier algebras.

\section{Stable balancing and the two real labels}

This section proves Theorem~\ref{thm:stable} by using the geometric maps just constructed to control the Busby invariants.

First specify the local finite-place model. Let $V=\Z\cup\{\infty\}$, with isolated integers and neighbourhoods $\{\infty\}\cup\{m\geq M\}$ at $\infty$. For $\ell=\log p$, let $\Z$ act on $V\times\R$ by
\[
 n\cdot(m,s)=(m+n,s+n\ell),\qquad
 n\cdot(\infty,s)=(\infty,s+n\ell).
\]
The map $\Q_p/\Z_p^\times\to V$ given by the valuation, with $0\mapsto\infty$, is a homeomorphism: its singleton and tail inverse images are $p^m\Z_p^\times$ and $p^M\Z_p$. Define
\[
 E_p^{\rm val}=C_0(V\times\R)\rtimes\Z.
\]
Its generic ideal is $C_0(\R)\otimes\K(\ell^2(\Z))$, with coordinate $t=s-m\ell$. Its boundary quotient is $C_0(\R)\rtimes\Z$ for translation by $\ell$.

The action on $V\times\R$ is free and proper: the real coordinate bounds the possible integer translates between two compact sets. The quotient has generic coordinate $t$ and boundary $\R/\ell\Z$, with convergence
\[
 t\to-\infty,\qquad [t]_\ell\to[s]_\ell.
\]
It is the unit circle together with the spiral
\[
 t\longmapsto(1+e^t)\exp(2\pi it/\ell).
\]
By the compatible edge imprimitivity already explained, the boundary unit is computed by the scalar lift $(1+e^t)^{-1}$ on the spiral and $1$ on the circle. Its exponential winds $-1$. Thus the finite local boundary coefficient is $-1$, in the same generic orientation as throughout the paper.

For the real place the reflection algebra is
\[
 E_\infty^{\rm ref}=C_0(\R)\rtimes C_2
 \cong\{F\in C_0([0,\infty),M_2(\C)):F(0)\text{ is diagonal}\}.
\]
Indeed the orbit matrix of $f_0+f_1u$ at $r>0$ is
\[
 \begin{pmatrix}f_0(r)&f_1(r)\\f_1(-r)&f_0(-r)\end{pmatrix}.
\]
At $r=0$ it commutes with the exchange matrix; conjugation by the normalized Hadamard matrix diagonalizes it. Conversely a continuous matrix path with this endpoint condition reconstructs continuous $f_0,f_1$ on $\R$, which proves the claimed isomorphism. The quotient is $\C^2$. Each endpoint rank-one projection lifts by $(1+r)^{-1}$ times that projection, so each has boundary $-1$ in the coordinate $t=\log r$. This is the same reflection calculation used by \citet[Lemma 5.6, pp.~26--27]{CC}.

Now make the local stabilizations precise. The finite edge of $S_{\{p\}}$, where the real coordinate is nonzero, is $M_2\otimes E_p^{\rm val}$: separate the real sign, take its logarithmic magnitude, and cross the freely translated sign coordinate. Let $\widehat E_p$ be this edge tensored with $\K(\ell^2(L_{\{p\}})\otimes H_0)$. Let $\widehat E_\infty=\widetilde S_\varnothing$. Their generic ideals are all the common $J$ under Lemma~\ref{lem:scalarcomparison}. Denote their quotients by $\widehat B_v$. The finite quotient is stably isomorphic to $C(S^1)$ by free proper translation and stable Morita equivalence \citep[Theorems 1.1--1.2, pp.~350--352]{BGR}; the real quotient is $\C^2\otimes\K$. These observations justify the quotient descriptions in Theorem~\ref{thm:stable} without changing any ideal coordinate.

\begin{proof}[Proof of Theorem~\ref{thm:stable}]
Let $\tau_v:\widehat B_v\to Q(J)=M(J)/J$ be the Busby maps of the specified local extensions. Their connecting coefficients are those just computed. Choose isometries $S_v$ on $H_0$, indexed by all finite primes and $\infty$, with mutually orthogonal ranges; partitioning a countable basis into countably many infinite subsets gives such a family. Put
\[
 s_v=1\otimes1\otimes S_v\in M(J).
\]
For $\mathcal B_P=\bigoplus_{v\in P\cup\{\infty\}}\widehat B_v$ define
\[
 \tau_P^{\rm bal}((a_v))=
 \sum_v\pi(s_v)\tau_v(a_v)\pi(s_v)^*.
\]
The cross-products vanish by orthogonality, so this is a $*$-homomorphism. Its pullback
\[
 E_P^{\rm bal}=
 \{(a,m)\in \mathcal B_P\oplus M(J):\tau_P^{\rm bal}(a)=\pi(m)\}
\]
is an extension by $J$. This is the stable Busby sum \citep[\S\S15.3, 15.6]{Blackadar}. The direct sum of the local extensions maps into it by
\[
 ((a_v,m_v))\longmapsto
 \left((a_v),\sum_vs_vm_vs_v^*\right).
\]
Each ideal corner $b\mapsto s_vbs_v^*$ induces the identity on $K_1(J)$: it carries a rank-one projection to a rank-one projection, and these projections are joined by a unitary path in a finite-dimensional subspace of the infinite Hilbert space. Suspending that path gives the required equality of generators. Thus the direct-sum ideal map induces addition, and naturality gives the stated negative-sum row.

The middle algebra is stable as well. Its concrete realization in $M(J)$ is
\[
 E_P^{\rm bal}=J+\sum_{v\in P\cup\{\infty\}}s_v\widehat E_vs_v^*:
\]
choose local lifts of each quotient coordinate, and the remaining multiplier lies in $J$. Each $\widehat E_v$ has the common tensor factor $\K(H_0)$. Choose the ranges of all $S_v$ to partition $H_0$, and a unitary $U_0:H_0\otimes\ell^2(\N)\to H_0$. The block unitary $U$ determined by
\[
 U(S_v\xi\otimes\eta)=S_vU_0(\xi\otimes\eta)
\]
carries $J\otimes\K$ onto $J$ and each stabilized summand onto the same summand. It therefore gives $E_P^{\rm bal}\otimes\K\cong E_P^{\rm bal}$.

The fixed place-indexed isometries also give injective maps
\[
 E_P^{\rm bal}\longrightarrow E_{P'}^{\rm bal},
 \qquad(a,m)\longmapsto((a,0),m).
\]
They preserve the Busby relation and compose exactly.

It remains to prove the comparison with the natural extension at the multiplier level. For $p\in P$, restrict $\Psi^{\rm sc}_{\{p\},P}$ to the finite edge $\widehat E_p$ and compose with $\iota_P$. For $v=\infty$, use $\Psi^{\rm sc}_{\varnothing,P}$ followed by $\iota_P$. These are injective local extension maps
\[
 \eta_{v,P}:\widehat E_v\longrightarrow\widetilde N_P
\]
with the same ideal map $j_P(b)=\one_{K_P}\otimes b$. Let $\beta_{v,P}:\widehat B_v\to\widetilde D_{v,P}$ be their quotient maps, and let $\sigma_P$ be the natural Busby map into $Q(C(K_P)\otimes J)$.

Since $j_P$ is nondegenerate it extends to a multiplier map $\overline j_P$ and a corona map $\dot j_P$. For each local extension, multiplication on its ideal gives the exact identity
\begin{equation}\label{eq:localbusbyidentity}
 \sigma_P\beta_{v,P}=\dot j_P\tau_v.
\end{equation}
In detail, for a local lift $e$ and $b\in J$ one has $\eta_{v,P}(e)j_P(b)=j_P(eb)$. The span $j_P(J)(C(K_P)\otimes J)$ is dense in the target ideal, so its multiplier action is uniquely $\overline j_P$ applied to the multiplier of $e$. Passing to the corona proves \eqref{eq:localbusbyidentity}.

The isometries $1\otimes S_v$ on the final $H_0$ factor belong also to $M(\widetilde N_P)$ and induce multipliers on its boundary quotient. Define
\[
 \alpha_P((a_v))=\sum_v s_v\beta_{v,P}(a_v)s_v^*,
\]
using these target multipliers. This map is injective: compress by $s_v^*$ and $s_v$, and use injectivity of $\beta_{v,P}$. Equation \eqref{eq:localbusbyidentity} and the fixed matrix factor give
\begin{equation}\label{eq:globalbusbyidentity}
 \sigma_P\alpha_P=\dot j_P\tau_P^{\rm bal}.
\end{equation}
The natural generic ideal is essential by Lemma~\ref{lem:scalarcomparison}, so its extension is its Busby pullback. Consequently
\[
 \Phi_P(a,m)=(\alpha_P(a),\overline j_P(m))
\]
is a well-defined injective homomorphism from $E_P^{\rm bal}$ to $\widetilde N_P$. Injectivity follows from injectivity of $\alpha_P$ and $\overline j_P$.

Proposition~\ref{prop:coherence} makes every local map $\eta_{v,P}$ compatible with addition of primes, and its matrix formula leaves the final $H_0$ factor unchanged. Hence the quotient maps $\alpha_P$ commute with the transition maps. On ideals both composite comparison maps are $j_{P'}$. They therefore agree on middle algebras too: their multiplier actions are the same on the nondegenerate image of $J$, and the target generic ideal is essential. This proves all the claimed coherence.

Finally the boundary row is surjective and has kernel
\[
 \{x\in\Z^{|P|+2}:\textstyle\sum x_i=0\}.
\]
For a fixed ordering of the quotient coordinates, its basis is $e_i-e_{i+1}$: the coefficient of that vector in a sum-zero $x$ is $x_1+\cdots+x_i$. Their Gram matrix has $2$ on the diagonal, $-1$ on adjacent entries, and zero elsewhere. It is the type-$A$ Cartan matrix. Since $K_0(J)=0$ and $K_1(\mathcal B_P)=\Z^{|P|}$, six-term exactness gives the two asserted middle $K$-groups.
\end{proof}

The geometric scalar algebra $\widetilde S_P$ itself has the same negative-sum boundary row and hence the same two $K$-groups. For $P\ne\varnothing$, this also follows directly by naturality with $\iota_P$: the map $K_1(J)=\Z\to C(K_P,\Z)$ is injective, finite quotient units pull back to constant units, and each real character pulls back to a rank-one class. For $P=\varnothing$ it is the reflection calculation above. The case of two finite primes is the result of \citet[\S5.4]{CC}. No isomorphism between $\widetilde S_P$ and $E_P^{\rm bal}$ is asserted; the explicit local maps suffice for \eqref{eq:globalbusbyidentity}.

\begin{proposition}[Comparison on the distinguished classes]\label{prop:labels}
For nonempty $P$, send each finite unit class to the constant class in $K_0(D_{p,P})$, and send both real character classes to the common class of $e_+$ and $e_-$ in $K_0(D_{\infty,P})$. Together with the inclusion of constant functions $\Z\to C(K_P,\Z)$, these maps give a commuting square between the stable scalar boundary row and \eqref{eq:delta}. The kernel of the quotient-label map is exactly $\Z(e_+-e_-)$. Its image has a balancing kernel of rank $|P|$.
\end{proposition}
\begin{proof}
Each finite unit lies in its own direct summand, where the constant function is nonzero and of infinite order. The real unit class has the same properties in its summand, but the two named projections have equal class by Proposition~\ref{prop:realboundary}. No further relation can mix distinct direct summands. This proves the label-map kernel. Every named generator has boundary $-\one_{K_P}$, so the square commutes. After identifying the two real labels, the negative-sum map has $|P|+1$ independent source coordinates and therefore kernel rank $|P|$.
\end{proof}

In particular a root lattice obtained from the scalar row is not the full $K_0(N_P)$. The three-prime example gives rank seven for the latter, while the distinguished constant-span kernel has rank three. Likewise, the ordinary augmentation of a free group on $r$ labelled sectors has kernel $A_{r-1}$ by the partial-sum calculation above. That calculation belongs to a chosen free linearization. Additive group completion of an idempotent monoid itself is zero: $x+x=x$ implies $[x]=0$ in any group. This prevents interpreting the augmentation lattice as the additive group completion of Boolean support data.

\section{Questions left by the calculation}

The results above settle stable existence of the place-balancing extension and compute the full connecting map of the semilocal one-zero extension. The following questions delimit what the calculation determines.

Theorem~\ref{thm:presentation} determines the extension class in \eqref{eq:mainK1} by explicit generators, relations, and finite transition maps. A further question is to classify its middle group as a direct sum of familiar abelian groups for arbitrary $P$, or to characterize exactly when that sequence splits. The quotient is a direct sum of localized-integer groups and need not be free. For $P=\{2,3,5\}$ it is
\[
 (\tfrac14\Z[1/15])^2\oplus\Z[1/10]^8\oplus\Z[1/6]^4.
\]
The precise remaining problem is whether the relations \eqref{eq:coboundarylift} can be removed by changing the lifts $\kappa_p$ by homomorphisms to $\coker\delta_P$. Equivalently, is the explicitly presented class zero in the corresponding $\operatorname{Ext}^1_{\Z}$ group? A boundary row and its finite Smith form do not answer this splitting question.

Theorem~\ref{thm:stable} settles the embedding and finite-set compatibility problem for the specified local models. A classification of every extension with the same integer row would be a different problem: the row records only one part of its Busby invariant. We make no such classification claim.

Finally, a geometric relation to modules over zero-adjoined rings requires a defined functor, not only a shared use of support sets. The semilocal crossed product proved here is an arithmetic construction in the sense of Connes, Consani and Marcolli. A link to inverse-semiring geometry remains a question.

\section*{Note on method and provenance}

The project began with a conversation in the Gemini app with Gemini 2.5 Pro. Gemini 2.5 Deep Think, in Deep Think mode, wrote the first papers in November 2025; versions were deposited on Zenodo under the name HI-AI from 6 November 2025 and later withdrawn. A GPT-5-series Pro model in ChatGPT Pro mode, at reasoning effort ultra, extracted algebraic material in March 2026. Further unpublished notes were written by GPT-5-series Pro models in that mode and effort during April--July 2026; the operator-algebra workbench was developed by Codex runs at reasoning effort ultra during August--September 2026. The November 2026 comparison notes were written by Claude Opus 5.5 at effort max, with consultations by Claude Fable 5.2 at effort max and Claude Opus 5.5 at effort high. That work did not know the earlier paper of Faul, Goswami, Joubert and Manuell.

The present mathematics and text were produced by GPT-6 Astra Pro. The current effort setting is not exposed. Statement-first re-proofs by the same model, proof comparisons, and independent finite computations are recorded separately. No human has checked the paper. Our text is dedicated to the public domain under CC0 1.0. Programs and mathematical records belong at \url{\workbenchurl}; the permanent record directory is \url{\recordurl}.

\begin{thebibliography}{9}
\bibitem[Blackadar(1998)]{Blackadar}
Blackadar, B. (1998). \emph{$K$-theory for operator algebras} (2nd ed.). Cambridge University Press. Author's electronic text: \url{https://www.bruceblackadar.com/Mathematics/book6.pdf}.

\bibitem[Brown, Green, and Rieffel(1977)]{BGR}
Brown, L. G., Green, P., \& Rieffel, M. A. (1977). Stable isomorphism and strong Morita equivalence of $C^*$-algebras. \emph{Pacific Journal of Mathematics, 71}(2), 349--363. \url{https://doi.org/10.2140/pjm.1977.71.349}.

\bibitem[Connes and Consani(2025)]{CC}
Connes, A., \& Consani, C. (2025). \emph{Knots, primes and class field theory}. arXiv:2501.06560v1. \url{https://doi.org/10.48550/arXiv.2501.06560}.

\bibitem[Connes, Consani, and Marcolli(2007)]{CCM}
Connes, A., Consani, C., \& Marcolli, M. (2007). \emph{The Weil proof and the geometry of the adeles class space}. arXiv:math/0703392v1. \url{https://arxiv.org/abs/math/0703392}.

\bibitem[Faul, Goswami, Joubert, and Manuell(2024)]{FGJM}
Faul, P. F., Goswami, A., Joubert, G., \& Manuell, G. (2024). \emph{The case for inverse semirings}. arXiv:2411.12318. \url{https://arxiv.org/abs/2411.12318}.

\bibitem[Guo, Proietti, and Wang(2024)]{GPW}
Guo, H., Proietti, V., \& Wang, H. (2024). A geometric Elliott invariant and noncommutative rigidity of mapping tori. \emph{Journal of Functional Analysis, 287}(11), 110625. \url{https://doi.org/10.1016/j.jfa.2024.110625}. arXiv version: \url{https://arxiv.org/abs/2207.06840v2}.

\bibitem[Green(1977)]{Green}
Green, P. (1977). $C^*$-algebras of transformation groups with smooth orbit space. \emph{Pacific Journal of Mathematics, 72}(1), 71--97. \url{https://doi.org/10.2140/pjm.1977.72.71}.
\end{thebibliography}
\end{document}
